%%%%%%%%%%%%
%
% $Autor: Wings $
% $Datum: 2019-03-05 08:03:15Z $
% $Pfad: TemplateSensor $
% $Version: 4250 $
% !TeX spellcheck = en_GB/de_DE
% !TeX encoding = utf8
% !TeX root = filename 
% !TeX TXS-program:bibliography = txs:///biber
%
%%%%%%%%%%%%

% Structure
\chapter{\textcolor{red}{Lichtschranke}}
 
\section{General}

Der Begriff Sensor stammt aus dem Lateinischen und bedeutet sinngemäß \glqq Fühler \grqq \space. 
Sensoren dienen zur qualitativen und quantitaiven Auffassung von physikalischen, chemischen, biologischen, klimatischen und medizinischen Größen . 
Hierbei bestehen Sensoren immer aus einem Sensorelement, welches sich mit Veränderung der physikalischen Gegebenheiten ebenfalls verändert. 
Besonders die Veränderung der elektrischen Eigenschaften des Sensorelementes ist hierbei von Interesse. Ein Beispiel sind Materialien, welche Ihren Widerstand bei Erwärmung verändern.
Zu jedem Sensorelement wird zudem eine Auswerteelektronik benötigt. Die Auswerteelektronik wandelt dann die beispielhafte Veränderung der elektrischen Eigenschaften in ein eindeutiges analoges oder digitales Signal. Dieses Signal kann dann das Verabreitende System, beispielweise ein Mikrocontroller übergeben werden \cite{Hering:2023}. 


cite books

\section{Specific Sensor/Actor}

cite board

\section{Specification}

\begin{itemize}
  \item Cite the data sheet
  \item Extract te information from the data sheet
  \item Circuit Diagram
\end{itemize}

\section{Library}

\subsection{Description}

\subsection{Installation}

\begin{itemize}
    \item data types
    \item data structure
    \item data quantity
    \item data quality
\end{itemize}

\subsection{Functions}

\section{Example}

In this section, a simple example is described in detail, which demonstrates how the library supports the sensor/actor.


\subsection{Manual}

\subsection{Inside the Example}

\subsection{Code}

\subsection{Files}



\section{Calibration}

cite method

\section{Simple Code}


\section{Simple Application}



\section{Tests}

\subsection{Simple Function Test}

\subsection{Test all Functions}



\section{Further Readings}

